\subsection{\texorpdfstring{The spaces \(\mathcal{F}_{\mathrm{F}}^{p}\)}{The spaces F sub F to the p}}

Let F be a Banach space, \(\mathcal{F}(X;F)\) (or simply \(F_{F}\) ) the vector space of all mappings of X into F. For \(1 \leqslant p < +\infty\) we will denote by \(\mathcal{F}^{p}(X,\mu;F)\) or \(\mathcal{F}_{\mathrm{F}}^{p}(X,\mu)\) , or simply \(\mathcal{F}_{\mathrm{F}}^{p}(\mu)\) , or \(\mathcal{F}_{\mathrm{F}}^{p}\) (if no confusion can result), the set of mappings f of X into F such that \(N_{p}(\mathbf{f}) < +\infty\) (we write \(F^{p}\) instead of \(F_{R}^{p}\) ). It is clear that \(\mathcal{F}_{\mathrm{F}}^{p}(|\mu|) = \mathcal{F}_{\mathrm{F}}^{p}(\mu)\) . It follows at once from Prop. 2 of No.~2 that \(F_{F}^{p}\) is a linear subspace of \(F_{F}\) and that \(N_{p}(\mathbf{f})\) is a semi-norm on this space. We shall always assume (absent express mention to the contrary) that \(F_{F}^{p}\) is equipped with the topology defined by this semi-norm; we shall say that this topology is the topology of convergence in mean of order p (for p = 1, we call it simply the topology of convergence in mean; for p = 2, one also says «topology of convergence in the quadratic mean»). We shall say that a filter G on \(F_{F}^{p}\) (resp. a sequence \((\mathbf{f}_{n})\) of elements of \(F_{F}^{p}\) ) that converges to f for this topology converges in mean of order p to f; this means, therefore, that \(N_{p}(\mathbf{g}-\mathbf{f})\) tends to 0 with respect to G (resp. that \(N_{p}(\mathbf{f}_{n}-\mathbf{f})\) tends to 0 as n tends to infinity).

This terminology extends at once to the case that the functions \(f_{n}\) and the function f are only defined almost everywhere (or have values in \(\overline{R}\) , and are defined and finite almost everywhere).

Note that the locally convex space \(F_{F}^{p}\) is in general not Hausdorff; the closure of 0 in this space is the linear subspace \(N_{F}\) of negligible mappings of X into F (No.~2, Prop. 3).

Remark. --- Let F be a Banach space over the field C of complex numbers; then, for every function \(f \in F_{F}^{p}\) and every complex number \(\alpha\) , \(\alpha f\) belongs to \(F_{F}^{p}\) and \(N_{p}(\alpha f) = |\alpha| N_{p}(f)\) ; in other words, \(F_{F}^{p}\) is also a vector space over C, and \(N_{p}(f)\) is a semi-norm on this complex vector space (TVS, II, §1).

PROPOSITION 4. --- Let B be a filter base on \(F_{F}^{p}\) . Assume that there exists a compact set \(K \subset X\) such that, for every set \(M \in B\) , all of the mappings \(f \in M\) have their support contained in K. Under these conditions, if B converges uniformly in X to \(f_{0}\) , then \(f_{0}\) belongs to \(F_{F}^{p}\) and B converges in mean of order p to \(f_{0}\) .

It comes to the same thing to say that, on the set of mappings \(f \in F_{F}^{p}\) whose support is contained in a fixed compact set, the topology of uniform convergence is finer than the topology of convergence in mean of order p.

For, let \(h\) be a continuous mapping of \(X\) into \([0,1]\) , with compact support, equal to 1 on \(K\) (Ch. III, §1, No.~2, Lemma 1). For every \(\varepsilon > 0\) , there exists an \(M \in \mathfrak{B}\) such that, for every mapping \(\mathbf{f} \in M\) , \(|\mathbf{f}(x) - \mathbf{f}_0(x)| \leqslant \varepsilon h(x)\) for all \(x \in X\) . From this, it follows that \(N sub p(\mathbf{f} - \mathbf{f}_0) \leqslant \varepsilon N sub p(h)\) , whence the proposition.

PROPOSITION 5. --- The locally convex space \(\mathcal{F}_{\mathrm{F}}^{p}\) is complete.

Since the Hausdorff space associated with \(\mathcal{F}_{\mathrm{F}}^{p}\) is a normed space, it suffices to prove that every Cauchy sequence \((\mathbf{f}_n)\) in \(\mathcal{F}_{\mathrm{F}}^{p}\) has a limit for the topology of convergence in mean of order \(p\) (GT, IX, §2, No.~6, Prop. 9). By hypothesis, for every \(\varepsilon > 0\) there exists an integer \(m_0\) such that the relations \(m \geqslant m_0\) , \(n \geqslant m_0\) imply \(\mathrm{N}_p(\mathbf{f}_n - \mathbf{f}_m) \leqslant \varepsilon\) . One may therefore define, by induction on \(k\) , a strictly increasing sequence \((n_k)\) of integers \(\geqslant 0\) such that \(\mathrm{N}_p(\mathbf{f}_{n_{k+1}} - \mathbf{f}_{n_k}) \leqslant 2^{-k}\) . If we show that the series with general term \(\mathbf{g}_k = \mathbf{f}_{n_{k+1}} - \mathbf{f}_{n_k}\) ( \(k \geqslant 1\) ) is convergent in mean of order \(p\) , then it will have a sum \(\mathbf{g} \in \mathcal{F}_{\mathrm{F}}^p\) , and \(\mathbf{f} = \mathbf{g} + \mathbf{f}_{n_1}\) will be the limit of the sequence \((\mathbf{f}_{n_k})\) in \(\mathcal{F}_{\mathrm{F}}^p\) ; \(\mathbf{f}\) will then be a cluster point of the sequence \((\mathbf{f}_n)\) ; since this sequence is a Cauchy sequence, it will have f as limit, and Prop. 5 will have been proved (GT, II, §3, No.~2, Cor. 2 of Prop. 5).

Prop. 5 is thus a consequence of the following proposition:

PROPOSITION 6. --- Let \((\mathbf{f}_{n})\) be a sequence of functions in \(F_{F}^{p}\) such that \(\sum_{n=1}^{\infty}\mathrm{N}_{p}(\mathbf{f}_{n})<+\infty\) . Under these conditions, the series with general term \(\mathbf{f}_{n}(x)\in\mathbf{F}\) is absolutely convergent almost everywhere in X. Setting \(\mathbf{f}(x)=\sum_{n=1}^{\infty}\mathbf{f}_{n}(x)\) at the points where the series converges, and \(\mathbf{f}(x)=0\) elsewhere, the function f belongs to \(F_{F}^{p}\) and is the sum of the series with general term \(f_{n}\) (for the topology of convergence in mean of order p); more precisely, for every \(n\geqslant0\) ,

\[
\mathrm{N} _ {p} \left(\mathbf {f} - \sum_ {k = 1} ^ {n} \mathbf {f} _ {k}\right) \leqslant \sum_ {k = n + 1} ^ {\infty} \mathrm{N} _ {p} (\mathbf {f} _ {k}).\tag{8}
\]

Consider the positive function (finite or not) \(g(x) = \sum_{n=1}^{\infty} |\mathbf{f}_n(x)|\) . By the countable convexity theorem (No.~2, Th. 1),

\[
\mathrm{N} _ {p} (g) \leqslant \sum_ {n = 1} ^ {\infty} \mathrm{N} _ {p} (\mathbf {f} _ {n}) <   + \infty ;
\]

thus \(g\) is finite almost everywhere (§2, No.~3, Prop. 7), which means that the series with general term \(\mathbf{f}_n(x)\) is absolutely convergent almost everywhere. Since \(F\) is complete, this series is convergent almost everywhere and, for every \(x \in X\) , \(|\mathbf{f}(x)| \leqslant \sum_{n=1}^{\infty} |\mathbf{f}_n(x)| = g(x)\) , whence

\[
\mathrm{N} _ {p} (\mathbf {f}) \leqslant \mathrm{N} _ {p} (g) \leqslant \sum_ {n = 1} ^ {\infty} \mathrm{N} _ {p} (\mathbf {f} _ {n}) <   + \infty ,
\]

which proves that \(\mathbf{f}\) belongs to \(\mathcal{F}_{\mathrm{F}}^{p}\) . On the other hand, for every integer \(n\) ,

\[
\left| \mathbf {f} (x) - \sum_ {k = 1} ^ {n} \mathbf {f} _ {k} (x) \right| \leqslant \sum_ {k = n + 1} ^ {\infty} | \mathbf {f} _ {k} (x) |
\]

almost everywhere, whence \(\mathrm{N}_{p}\left(\mathbf{f}-\sum_{k=1}^{n}\mathbf{f}_{k}\right)\leqslant\sum_{k=n+1}^{\infty}\mathrm{N}_{p}(\mathbf{f}_{k})\) . By hypothesis, the series with general term \(\mathrm{N}_{p}(\mathbf{f}_{n})\) is convergent; therefore, for every \(\varepsilon > 0\) , there exists an integer \(n\) such that \(\sum_{k=n+1}^{\infty} \mathrm{N}_{p}(\mathbf{f}_{k}) \leqslant \varepsilon\) , and the inequality (8) proves that \(\mathbf{f}\) is the sum of the series with general term \(\mathbf{f}_{n}\) , for the topology of convergence in mean of order \(p\) .

Propositions 5 and 6 are thus completely proved.

